\section{Related Work}
\label{sec:related}

\paragraph{Coding agents beyond software development.}
SWE-agent studies interfaces for repository navigation, code editing, and execution
\citep{yang2024sweagent}, while OpenHands provides a platform for agents that act through
code, a command line, and a browser \citep{wang2025openhands}. SWE-bench evaluates issue
resolution in real software repositories \citep{jimenez2024swebench}. Beyond software,
GAIA and AgentBench assess information gathering, tool use, and interaction with diverse
environments \citep{mialon2024gaia,liu2024agentbench}; ScienceAgentBench and MLE-bench
evaluate scientific data analysis and machine-learning engineering
\citep{chen2025scienceagentbench,chan2025mlebench}. These studies establish that agent
evaluation extends beyond code generation. Our focus is the transfer of a general-purpose
coding agent to tasks whose target outputs are non-coding, compared with domain specialists
using both published scores and executions on a shared backbone.

\paragraph{Domain-specific scaffolding.}
AutoGen supports multi-agent conversations \citep{wu2024autogen}; DSPy and AFlow optimize
language-model pipelines and agentic workflows \citep{khattab2024dspy,zhang2025aflow}.
Domain-focused systems instantiate this structure through medical expert discussions in
MedAgents \citep{tang2024medagents} and chemistry tool libraries in ChemCrow
\citep{bran2024chemcrow}. Additional structure also creates failure points: MAST organizes
observed multi-agent failures into system design, inter-agent misalignment, and task
verification categories \citep{cemri2025mast}. Building on these perspectives, we examine
whether longer workflows improve accuracy, whether decomposition preserves necessary
information, and when explicit domain conventions, task conditions, or verification help.
Our paired comparisons distinguish published performance from performance on our backbone;
trace examples motivate mechanisms whose causal effects require targeted ablations.

\paragraph{Reliable benchmark evaluation.}
GSM1k probes overfitting to a widely used reasoning benchmark \citep{zhang2024gsm1k}, and
prompt-format experiments show that meaning-preserving changes can alter measured performance
\citep{sclar2024format}. HELM standardizes evaluation conditions across scenarios and metrics
\citep{liang2023helm}, while the Language Model Evaluation Harness provides reusable task
and scoring implementations \citep{gao2023harness}. The NeurIPS reproducibility program
complements such infrastructure with code, checklists, and independent reproduction
\citep{pineau2021reproducibility}. Our audit connects each published score to its reference
implementation and the executable task used for comparison, checking instances, metrics,
answer keys, judges, and resources. These checks establish the conditions for interpreting
performance gaps; we treat them separately from findings about agent capabilities.

\paragraph{AI-assisted evidence synthesis.}
ASReview demonstrates how active learning can prioritize literature screening while retaining
a transparent review process \citep{vandeschoot2021asreview}. Our \emph{meta-analysis pipeline}
extends evidence collection to checking reported benchmark results and executing comparisons.
It does not estimate a pooled effect across heterogeneous metrics, and it is not an exhaustive
systematic review: the search was not preregistered and screening was not performed by two
independent human reviewers. Source records, adversarial checks, and human adjudication make
the process reviewable; an independent accuracy evaluation of the SOTA finder remains
unfinished (Section~\ref{sec:rigor}).
